\documentclass[11pt]{article}

\usepackage[T1]{fontenc}
\usepackage[utf8]{inputenc}
\usepackage{lmodern}
\usepackage{amsmath}
\usepackage{amssymb}
\usepackage{amsthm}
\usepackage{booktabs}
\usepackage{tabularx}
\usepackage{microtype}
\usepackage[margin=1in]{geometry}
\usepackage[hidelinks]{hyperref}

\DeclareMathOperator{\opt}{opt}
\DeclareMathOperator{\tree}{tree}
\DeclareMathOperator{\gapop}{gap}
\newcommand{\optx}{\opt_{\mathrm{XAG}}}
\newcommand{\treex}{\tree_{\mathrm{XAG}}}
\newcommand{\gapx}{\gapop_{\mathrm{XAG}}}
\newcommand{\treemig}{\tree_{\mathrm{MIG}}}

\theoremstyle{plain}
\newtheorem*{claim}{Claim}

\newcolumntype{Y}{>{\raggedright\arraybackslash}X}

\title{The Unit Gap Is Base-Dependent:\\
A Mechanically Verified Refutation and a Cross-Base Census}
\author{%
  Luiz Antonio Busnello \\
  Curious Tech Entrepreneur \\
  S\~{a}o Paulo, Brazil \\
  \texttt{lab@generantis.com.br} \\
}
\date{}

\begin{document}
\maketitle

\begin{abstract}
Krinkin's \emph{The Unit Gap: How Sharing Works in Boolean Circuits} (arXiv:2603.08033v2) asserts
that for every Boolean function the minimum formula size exceeds the minimum circuit size by at most
one gate --- Theorem~2: $\gapop \in \{0,1\}$ in the And-Inverter Graph (AIG) cost model --- with a
decomposition corollary bounding a shared-gate term $s \in \{0,1\}$ (Corollary~6). We refute both.
Under the paper's own definition of a formula (fan-out one at every gate), the parity of three
variables ($\oplus_3$) is a counterexample: $\opt(\oplus_3) = 6$ --- UNSAT for every gate count
$k = 1,\dots,5$, each certified by a DRAT proof checked with \texttt{drat-trim} --- while
$\tree(\oplus_3) = 9$, giving gap~3; in the optimal six-gate circuit the two output-child sub-DAGs
intersect in a three-gate parity-2 block, so Krinkin's own term $s = |D_a \cap D_b| = 3$.
\textbf{The refutation itself is classical.} The AIG basis with free inversions is $U_2$, where
Schnorr's theorem gives the circuit complexity of $\oplus_n$ as exactly $3n-3$, so
$\opt(\oplus_3) = 6$; Khrapchenko forces $\tree(\oplus_3) \ge 9$ leaves, hence $\ge 8$ gates (a
formula with $L$ leaves has $L-1$ gates). \textbf{So $\gapop(\oplus_3) \ge 2 > 1$ follows from
two published results and the basis identification above, with no computation of ours in the
chain: Theorem~2 was already contradicted by the literature.} The exact value 3 needs
$\tree(\oplus_3) = 9$, which is not in the refereed literature; it is certified here by the DRAT
chain, two independent dynamic programs and the Lean development, and agrees with an exhaustive
enumeration published informally in 2011. At $n = 4$ the exact gap is classical again: Tarui
determined the De~Morgan formula size of $\oplus_n$ for $n = 2^k$ to be exactly $n^2$ leaves,
giving $\gapop(\oplus_4) = 15 - 9 = 6$ --- the maximum of the census below. The error is locatable in the
paper: \S2 displays $\tree(f) = \min(1 + \opt(a) + \opt(b))$ (the DAG measure $\opt$ in the
children), while \S3's Bellman operator is $(Tv)(f) = \min(1 + v(a) + v(b))$ (the recursive
measure); these are incompatible characterizations of formula size, the first making Theorem~2 a
one-line consequence and the second making it false (its fixed point overshoots $\opt(\oplus_3)$ by
3). An exhaustive census over the 222 NPN classes of $n = 4$ shows the failure is structural: 72 of
222 classes have $\gapop \ge 2$, the maximum gap 6 attained by three distinct
NPN classes (distinguished by their XAG optima 6/5/3), one of which is parity-4. Finally, the
phenomenon is \textbf{base-dependent}. In the XOR-AND-graph (XAG) basis the unit-gap property holds
for $n \le 4$ (under a separately gated exact-synthesis pipeline whose UNSAT side is not
proof-certified): $\gapx = 0$ for all 256 functions of $n = 3$ and $\gapx \in \{0,1\}$ for all 222
classes of $n = 4$; a directed search over 25{,}373 distinct essential functions of $n = 5$
\textbf{found no separator}: all 25{,}169 gap-decided functions have $\gapx \in \{0,1\}$, and 204
remain unresolved. A replicated MIG census matching Soeken et al.'s published size
distribution bucket by bucket gives an aggregate external cross-check. The large AIG gaps are
consistent with the cost of simulating linear (parity) structure without a native XOR gate.
\end{abstract}

\section{Introduction}

The relationship between \textbf{formula size} (the smallest fan-out-one circuit for a function) and
\textbf{circuit size} (the smallest circuit of any fan-out) is one of the oldest quantitative
questions in Boolean complexity. Their difference measures how much \emph{sharing} --- reusing an
intermediate result at more than one place --- can buy. Krinkin's \emph{The Unit Gap} claims that in
the And-Inverter Graph model this difference is \textbf{at most one gate} for every Boolean function
(Theorem~2), with a decomposition corollary bounding the sharing (Corollary~6).

It is not true. This paper does three things.

\begin{enumerate}
  \item \textbf{We refute Theorem~2 and Corollary~6} at $n = 3$ (parity), where the gap is 3 and
  Krinkin's shared-gate term is 3. \emph{The counterexample is not new mathematics}: the
  contradiction with Theorem~2 was already implied by the published literature, since Schnorr
  fixes the circuit side and Khrapchenko bounds the formula side, and the two together give
  $\gapop(\oplus_3) \ge 2 > 1$. What we add is to
  notice it, to certify it mechanically (a DRAT chain at every $k = 1,\dots,5$ checked by an
  independent tool, two dynamic programs, and a Lean~4 formalization), and above all \textbf{to
  locate the error inside the paper}: \S2 and \S3 give incompatible characterizations of
  $\tree(f)$, one making Theorem~2 trivially true and the other making it false. A refutation that
  only exhibits a counterexample leaves the author guessing where the argument broke; this one does
  not.

  \item \textbf{We show the failure is structural}, computing the complete gap census over all 222
  NPN classes of $n = 4$: 32.4\% of classes violate the unit-gap bound, the maximum gap, 6, attained by
  three distinct NPN classes, one being parity.

  \item \textbf{We show the phenomenon is base-dependent.} The property, false in AIG, holds in the
  XAG basis for $n \le 4$ (exhaustive); a large directed search at $n = 5$ found no counterexample
  among the functions it decided. The large parity
  AIG gaps are consistent with the price of simulating parity without a native XOR gate.
\end{enumerate}

Our method is \textbf{exact synthesis via SAT}: ``is there a $k$-gate circuit for $f$?'' is a CNF
decided by a modern solver. The encoder is a normalized exact-size encoding (every non-output gate
used, no duplicate gates); by a minimality lemma, $\opt(f) = k$ implies the $k$-encoding is
satisfiable, so a lower bound $\opt(f) \ge m$ requires UNSAT at \emph{every} $k = 1,\dots,m-1$ ---
which is why our circuit lower bounds certify the whole chain, not just the last step. The package
was reviewed adversarially by four independent language-model families before external
communication; the mathematics rests on the certificates, not that review. We claim novelty only for
the refutation and for the gap-by-basis census; the individual optimal-size tables in XAG and MIG
are classical or previously published.

\section{Preliminaries}
\label{sec:prelim}

We quote the paper's definitions to fix the target precisely (arXiv:2603.08033v2, \S2;
notation-normalized).

\begin{quote}
\itshape
A circuit is an AIG where gates may have fan-out greater than one. A formula (or tree circuit) has
fan-out exactly one at every gate. We write $\tree(f)$ for the minimum formula size. Since every
formula is a circuit, $\tree(f) \ge \opt(f)$. The gap is $\gapop(f) = \tree(f) - \opt(f)$.
\end{quote}

And, for the decomposition (\S2):

\begin{quote}
\itshape
Every AIG gate computes $f = a \wedge b$ (with possible edge complementation). The minimum formula
size satisfies $\tree(f) = \min$ over $a,b$: $f = a \wedge b$ or $f = \neg(a \wedge b)$ of
$(1 + \opt(a) + \opt(b))$. The trivial decomposition $f = 1 \wedge f$ always achieves $1 + \opt(f)$.
\end{quote}

An \textbf{AIG} over $x_1,\dots,x_n$ and the constant 1 is a DAG of two-input AND gates with free
edge complementation; $\opt(f)$ is the minimum number of AND gates. With free complementation an
AND gate realises every binary Boolean function except XOR and XNOR at unit cost, so the AIG basis
is exactly $U_2$, and $\opt$ coincides gate for gate with $U_2$-circuit size. \emph{The
exclusion of XOR is what makes this basis the right one for the classical results used below:}
$U_2$ is the full binary basis $B_2$ \emph{minus} XOR and XNOR, so parity is not free in it. (The output gate has
fan-out zero; the ``fan-out one'' condition is Krinkin's, for the non-output gates of a tree.) An
\textbf{XAG} adds a native XOR gate; an \textbf{MIG} (Majority-Inverter Graph) uses 3-input majority gates. All exhaustive
results are for $n \le 4$.

\paragraph{Exact synthesis.} For fixed $(n, k)$ we one-hot--select each gate's operation and
operands among earlier nodes and constrain the output to compute $f$ up to polarity. The encoding is
\emph{normalized}: every non-output gate must be used and duplicate gate options are forbidden.
These are sound for minimality: SAT at $k$ yields a valid $k$-gate circuit, and if $\opt(f) = k$ a
minimality lemma guarantees the $k$-encoding is satisfiable. Hence $\opt(f)$ is the least $k$ with
SAT, and a lower bound $\opt(f) \ge m$ requires UNSAT at \emph{every} $k = 1,\dots,m-1$. Restricting
selection to fan-out one turns the same encoder into a $\tree(f)$ oracle (validated in
Appendix~\ref{app:encoders}).

\section{The refutation at \texorpdfstring{$n = 3$}{n = 3}}
\label{sec:refutation}

\paragraph{The paper gives incompatible characterizations of $\tree(f)$.} The \S2 identity puts the
DAG measure $\opt$ in the children; with $\opt$ there, the trivial decomposition $f = 1 \wedge f$
yields $\tree(f) \le 1 + \opt(f)$ for every $f$, so Theorem~2 is a one-line consequence --- but the
minimized object is not a tree (the children may be internally shared DAGs). In \S3, however, the
paper defines the Bellman operator (notation-normalized):

\begin{quote}
\itshape
$(Tv)(f) = \min$ over $a,b$: $f = a \wedge b$ of $(1 + v(a) + v(b))$; iterating $T$ from the base
case $v_0(x_i) = v_0(1) = 0$ computes $\tree(f)$.
\end{quote}

with the \textbf{recursive} measure $v$ in the children. At its fixed point this reads
$\tree(f) = \min(1 + \tree(a) + \tree(b))$, a different object from the \S2 identity: as global
characterizations of formula size the two coincide term-by-term only on children for which
$\opt = \tree$. Krinkin further states --- as a \emph{consequence of Theorem~2}, not a definition ---
that $T$'s fixed point ``overshoots $\opt(f)$ by exactly $\gapop(f) \in \{0,1\}$'', which collapses
together with Theorem~2. For $\oplus_3$ the recursive $T$ computes 9, overshooting $\opt = 6$ by 3;
under the \S3 reading, Theorem~2 is false.

\begin{itemize}
  \item \textbf{$\opt(\oplus_3) = 6$.} \emph{Upper bound:} the six-gate circuit
  \begin{equation}\label{eq:c6}
  \begin{aligned}
    g_1 &= x_1 \wedge \neg x_2, &\qquad g_3 &= \neg g_1 \wedge \neg g_2, &\qquad g_5 &= g_3 \wedge x_3,\\
    g_2 &= \neg x_1 \wedge x_2, &\qquad g_4 &= \neg g_3 \wedge \neg x_3, &\qquad g_6 &= \neg g_4 \wedge \neg g_5,
  \end{aligned}
  \end{equation}
  with output $\neg g_6$ (inversions are free, so the output polarity costs nothing). Here
  $g_3 = \neg(x_1 \oplus x_2)$ is the \emph{only} gate of fan-out two, which is exactly why the
  circuit is not a tree; the two children of the output, $g_4$ and $g_5$, have sub-DAGs meeting in
  $\{g_1, g_2, g_3\}$ --- this is the $s = 3$ of the Corollary~6 refutation below, and it can be
  read off the display. Verified by simulation on all eight inputs and in Lean~4.
  \emph{Lower bound:}
  UNSAT at \textbf{every $k = 1,2,3,4,5$} (kissat~4.0.4), each with a DRAT proof checked by
  \texttt{drat-trim} (\texttt{s VERIFIED}; transcript and hashes in Appendix~\ref{app:repro}). The
  full $k = 1,\dots,5$ chain --- not only $k = 5$ --- certifies $\opt \ge 6$ under the normalized
  exact-size encoding. \emph{The value is also classical:} since the AIG basis is $U_2$ (\S\ref{sec:prelim}),
  Schnorr's theorem --- the $U_2$ circuit complexity of $\oplus_n$ is exactly
  $3n-3$~\cite{schnorr1976}, an equality restated in~\cite{morizumi2015}\footnote{We cite the
  1976 original for the result and a modern, openly accessible source for the statement of the
  equality; the original was not consulted directly.} --- gives $\opt(\oplus_3) = 6$ and $\opt(\oplus_4) = 9$ directly, in
  agreement with the values computed here.

  \item \textbf{$\tree(\oplus_3) = 9$.} \emph{Upper bound:} the nine-gate formula is
  \eqref{eq:c6} with the shared block duplicated, which is what fan-out one forces:
  \begin{equation}\label{eq:w9}
    \neg\bigl(\,\neg P_1 \wedge \neg x_3\,\bigr) \wedge \neg\bigl(\,P_2 \wedge x_3\,\bigr),
    \qquad P_1 = P_2 = \neg(x_1 \oplus x_2),
  \end{equation}
  where each $P_i$ is a private three-gate copy of the block that \eqref{eq:c6} computes once as
  $g_3$. Counting: $3 + 3$ for the two copies, plus the three gates above them, is 9; the circuit
  pays $3$ for one copy plus the same three, which is 6. \textbf{The gap is the duplicated block},
  and its size is Krinkin's $s = 3$. Checked the same way as~\eqref{eq:c6}. \emph{Lower bound:} the exact fixed point of the tree recursion over all 256
  functions, independently confirmed by a layer-by-layer enumeration (new functions per cost
  $1,\dots,9$: 24, 64, 30, 80, 32, 0, 16, 0, 2 --- the two cost-9 functions are parity and its
  complement).\footnote{The gap distribution at $n = 3$ (\{gap 0: 214, 1: 40, 3: 2\}) is reproduced
  by \texttt{tree\_gap\_n3.py}; the per-cost histogram above was re-derived independently during
  review.}

  \item \textbf{Analytic fail-safe.} Khrapchenko~\cite{khrapchenko1971}: $L(\oplus_3) \ge |E|^2/(|A||B|) = 144/16 = 9$
  leaves, hence $\ge 8$ two-input gates. Together with the explicit six-gate circuit
  ($\opt \le 6$), $\gapop(\oplus_3) \ge 2$ follows \textbf{using only Khrapchenko and that circuit}
  --- independent of the tree DP and the DRAT chain. Substituting Schnorr's theorem for the
  explicit circuit makes the fail-safe entirely citational: $\opt(\oplus_3) = 6$~\cite{schnorr1976}
  and $\tree(\oplus_3) \ge 8$ (Khrapchenko) yield $\gapop(\oplus_3) \ge 2 > 1$ from two published
  classical results and the basis identification of \S\ref{sec:prelim}, with no computation of
  ours in the chain.

  \emph{Beyond the fail-safe, the exact value is also classical.} The identification of \S\ref{sec:prelim} extends from circuits
  to formulas: with free complementation every $U_2$ gate is an AND or an OR of its two arguments up
  to edge inversions, so pushing inversions to the leaves turns an AIG formula into a De~Morgan
  formula with the same number of binary gates and back. Hence
  $\tree = \tree_{U_2} = \tree_{\mathrm{De\,Morgan}}$, and a binary formula with $L$ leaves has
  $L - 1$ gates --- the same conversion already used for Khrapchenko above. So Khrapchenko gives
  $\tree(\oplus_3) \ge 9$ leaves, i.e. $\ge 8$ gates, and with $\opt(\oplus_3) = 6$ this already
  contradicts Theorem~2: $\gapop(\oplus_3) \ge 2 > 1$, from two published results and the
  identification above.

  The \emph{exact} value is a different matter. For $n = 2^k$ Tarui~\cite{tarui2010} showed that
  Khrapchenko's $n^2$ is attained, so the formula size is exactly $n^2$ leaves; this fixes
  $\tree(\oplus_4) = 16$ leaves, i.e. 15 gates, and hence $\gapop(\oplus_4) = 6$, the maximum gap
  found in the $n = 4$ census of \S\ref{sec:census} --- an external check on the census, not
  merely on the counterexample. For $n = 3$ we found no determination in the refereed
  literature.\footnote{The obvious candidate does not survive. Lee's rank
  technique~\cite{lee2007} was announced as determining the formula size of parity for every $n$,
  as $2^{\ell}(2^{\ell}+3r)$ leaves for $n = 2^{\ell}+r$, which would give 10 leaves at $n = 3$.
  Lee has since retracted that claim: the error is in his Proposition~2, which assumes an optimal
  selection function takes each $s_i$ to be a power of two, and he points to
  Rychkov~\cite{rychkov1994} for the current best lower bounds --- $n^2+3$ for odd $n \ge 5$ and
  $n^2+2$ for even $n \ge 6$ that are not powers of two, neither of which covers $n = 3$. The
  rank technique itself, and Lee's observation that no bound above $n^2$ is reachable by
  Khrapchenko's method, are unaffected.} Accordingly $\tree(\oplus_3) = 9$ is certified here
  rather than cited: by the tree dynamic program, independently re-implemented, and by the Lean
  development. The value is not new --- an exhaustive enumeration by Cox and Healy~\cite{cox2011}
  gives 9 AND/OR operators for $\oplus_3$ and 15 for $\oplus_4$, and writes out the optimal
  $n = 3$ formula --- and the agreement is an external check on both.
\end{itemize}

Hence $\gapop(\oplus_3) = 3$, refuting Theorem~2.

\paragraph{Corollary 6 is refuted with Krinkin's own definition.} He states (\S5;
notation-normalized):

\begin{quote}
\itshape
Let $C$ be an optimal AIG, and let $g$ be any gate computing $f = a \wedge b$. Let $D_a, D_b$ be the
sub-DAGs computing $a$ and $b$. Then $\opt(f) = 1 + \opt(a) + \opt(b) - s(a,b)$, where
$s(a,b) = |D_a \cap D_b| \in \{0,1\}$ counts the shared gates.
\end{quote}

Structurally, take $g$ to be the output gate of the six-gate optimum for $\oplus_3$: its children
$a, b$ share the three-gate parity-2 block (computing $\neg(x_1 \oplus x_2)$ in the Lean witness
\texttt{c6}; $x_1 \oplus x_2$ up to free output polarity), so $|D_a \cap D_b| = 3$ and, since each
child has $\opt = 4$, $s = 1 + 4 + 4 - 6 = 3 \notin \{0,1\}$. One child,
$(x_1 \oplus x_2) \wedge \neg x_3$, has its $\opt = 4$ certified by a DRAT chain (UNSAT at
$k = 1,2,3$; Appendix~\ref{app:repro}); the other child is NPN-equivalent (input/output polarity),
hence also $\opt = 4$. Krinkin's own $s \le 1$ derivation is separately unsound: he needs
$s \le \tree(f) - \opt(f)$, but the \S2 identity gives $\tree(f) \le 1 + \opt(a) + \opt(b)$, i.e.
$s \ge \tree(f) - \opt(f)$ --- the reverse of the comparison the proof requires. (For parity itself
$s = \tree - \opt = 3$, an equality, so Krinkin's first inequality happens to hold there; his chain
then breaks at the \emph{second} step, $\tree - \opt \le 1$, which is Theorem~2.) Krinkin's Table~3
reports $s = 1$ for the complete $\opt$-6 $n = 3$ bucket, which includes the parity decompositions
--- consistent with his having imposed the bound, and in conflict with the structural count
$s = 3$.

\paragraph{Fate of the other results.} Theorem~7 (a classification assuming $\gapop = 1$) loses its
proof, which depends on Corollary~6; its conditional statement under the standard definition is not
refuted by $\oplus_3$ (gap~3) and is left unproven. \textbf{The statements of Theorems 3 and 4 stand.} Theorem~3's
published proof applies Lemma~1 (gate elimination) to the sub-DAG $S$ below the shared gate $g$ and
writes $|S| \ge k-1$, which fails when $g$ is at input level; a corrected incidence-count proof is
in Appendix~\ref{app:thm3}.

\paragraph{Diagnostic signature.} In Krinkin's complete $n = 3$ table (Table~1), the two functions
with $\gapop = 1$ at $\opt = 6$ are parity and its complement --- precisely what the \S2 identity
(with $\opt$ in the children: $1 + \opt(\text{parity}) + \opt(1) = 7$) reports, while the true
formula size is 9. The error is visible in the paper's data.

\paragraph{Verification.} The counterexample is formalized in Lean~4.31.0 without Mathlib
(\texttt{UnitGap.lean}): the witnesses and the DP's structural framework are kernel-checked
(\texttt{decide}, induction); the finite fact $\mathrm{par3} \notin D_8$ --- hence
$\tree(\oplus_3) \ge 9$ --- is discharged by \texttt{native\_decide} (compiler-trusted, not kernel).
We report this precisely rather than as blanket ``kernel-checked'' (see \S\ref{sec:verification}).

\section{The failure is structural: AIG gap census at \texorpdfstring{$n = 4$}{n = 4}}
\label{sec:census}

We computed $\tree(f)$ for all 65{,}536 functions by a layered dynamic program (independently
re-implemented as a global Bellman fixed point with explicit polarities --- all 65{,}536 cells
agree), joined with the complete exact $\opt$ catalog. The distribution of $\gapop$ over the 222 NPN
classes:

\begin{table}[h]
\centering
\begin{tabular}{lrrrrrrr}
\toprule
$\gapop$ & 0 & 1 & 2 & 3 & 4 & 5 & 6 \\
\midrule
classes & 93 & 57 & 40 & 13 & 14 & 2 & 3 \\
\bottomrule
\end{tabular}
\end{table}

\noindent
\textbf{72 of the 222 NPN classes (32.4\% of classes) have $\gapop \ge 2$} (a fraction of classes,
not of all 65{,}536 functions, whose orbit sizes differ). The maximum gap 6 is attained by three
\textbf{distinct} NPN classes \texttt{\{0x1668, 0x16e9, 0x6996\}} --- all with $\opt = 9$,
$\tree = 15$ --- distinguished by their XAG optima ($\optx = 6, 5, 3$ respectively; $\optx$ is
NPN-invariant, so distinct values force distinct classes). One of the three is parity-4
(\texttt{0x6996}, $\optx = 3$). For parity-4, Khrapchenko's bound ($\ge 16$ leaves $\implies \ge 15$
gates) meets $\tree = 15$ with equality. The embedded $n = 3$ functions reproduce the $n = 3$ table
256/256. Parity is thus \textbf{a} witness of the maximum gap, not its exclusive source: the census
is consistent with, but does not by itself isolate, ``linear structure without native XOR'' as the
mechanism for all 72 violating classes.

\section{Base dependence: XAG, MIG, and a directed search at \texorpdfstring{$n = 5$}{n = 5}}
\label{sec:basedep}

The large parity AIG gaps are consistent with the cost of building parity from AND gates, suggesting a
basis artifact.

\paragraph{XAG, $n \le 4$ (exhaustive).} With a native XOR gate, the pipeline (encoder validated by
an independent exhaustive enumeration at $n = 3$ and by a formula-encoder gate on all 222 $n = 4$
classes) gives $\optx$ maximum 7 (strictly below AIG in 190/222 classes; $\optx(\oplus_4) = 3$) and
$\treex$ maximum 7. The gap census (\texttt{census/xag-n4/npn4\_xag\_gap.csv}): $n = 3$ ---
$\gapx = 0$ for all 256 functions; $n = 4$ --- $\gapx \in \{0,1\}$, distributed $\{0: 218, 1: 4\}$.
The unit-gap property, false in AIG, \textbf{holds in XAG for $n \le 4$} under this separately gated
encoder whose UNSAT instances are not proof-certified (\S\ref{sec:verification}).

\paragraph{XAG, $n = 5$ (directed search).} Exhaustive census is infeasible ($2^{32}$ truth tables).
We searched for a \emph{separator} ($\gapx \ge 2$) by generating candidate functions from random XAG
circuits with sharing and XOR bias, keeping those essential in all five variables, and deciding
$\optx$ and $\treex$ exactly by SAT (ascending search, 20\,s cap). This is a
\textbf{generator-reachable, solver-decided sample}, not a probability statement: functions whose
$\opt$ timed out or exceeded the cap are not written and thus not counted, so the sample is
conditional on easy circuit optimization. The formula encoder was gated before the search (matching
the $n = 4$ DP on all 222 classes, including the four with $\gapop = 1$); an early false positive
(parity-5) was fixed by ascending from $\opt$ (Appendix~\ref{app:encoders}).

Over \textbf{25{,}373 distinct} essential functions of five variables (deduplicated across workers):

\begin{table}[h]
\centering
\begin{tabular}{lrr}
\toprule
verdict & functions & \% \\
\midrule
gap 0 & 24{,}875 & 98.04 \\
gap 1 & 294 & 1.16 \\
inconclusive (formula-search timeout) & 204 & 0.80 \\
\textbf{separator ($\gapx \ge 2$)} & \textbf{0} & \textbf{0} \\
\bottomrule
\end{tabular}
\end{table}

\noindent
\textbf{No separator was certified.} All \textbf{25{,}169 gap-decided} functions have
$\gapx \in \{0,1\}$; the 204 inconclusive cases are formula-search timeouts and could in principle
be separators. As an added check, all 294 gap-1 functions and a sample of 200 gap-0 functions were
re-decoded and simulation-verified (optimal circuit and formula both compute the truth table);
494/494 passed (Appendix~\ref{app:repro}). This is a \textbf{directed sample, not a proof}: a
separator of high $\opt$, or outside the generation bias, is not excluded. What the search
establishes is that no separator appears among 25{,}169 gap-decided, generator-reachable essential
functions of $n = 5$.

\paragraph{MIG (aggregate cross-check).} Soeken, Amarù, Gaillardon and De Micheli~\cite{soeken2016}
published the optimal MIG size distribution for the 222 classes of $n = 4$. We re-derived the census
with a different solver (kissat vs.\ Z3) and encoding and matched the published \textbf{size
distribution} $\{0{:}2, 1{:}2, 2{:}5, 3{:}18, 4{:}42, 5{:}117, 6{:}35, 7{:}1\}$ \textbf{bucket by
bucket}, including the unique 7-node class. (Their table publishes bucket counts, not a per-class
map; equal aggregate distributions cross-check the shared exact-synthesis machinery, not the
per-class assignments.)

\paragraph{Cross-base summary ($n = 4$).}

\begin{table}[h]
\centering
\begin{tabular}{lccc}
\toprule
 & AIG & XAG & MIG \\
\midrule
max $\opt$ & 10 & 7 & 7 \\
$\opt(\oplus_4) / \tree(\oplus_4)$ & 9 / 15 & 3 / 3 & 6 / --- \\
gap census & up to 6; 32.4\% of classes $\ge 2$ & $\{0,1\}$ & deferred \\
\bottomrule
\end{tabular}
\end{table}

\noindent
($\treemig$ requires a ternary DP whose naive form is cubic in set sizes; deferred.)

\section{Verification}
\label{sec:verification}

\begin{table}[h]
\centering
\begin{tabularx}{\textwidth}{@{}p{0.30\textwidth}Y@{}}
\toprule
Claim & Level of certification \\
\midrule
$\opt(\oplus_3) = 6$ (lower bound) &
DRAT proofs for UNSAT at \textbf{every $k = 1,\dots,5$}, \texttt{drat-trim} \texttt{s VERIFIED};
transcript + hashes archived (Appendix~\ref{app:repro}) \\
\addlinespace
$\opt$ of the shared child $= 4$; other child NPN-equivalent &
DRAT proofs for UNSAT at $k = 1,2,3$ on the representative child, \texttt{drat-trim}
\texttt{s VERIFIED}; the second child is NPN-equivalent \\
\addlinespace
$\tree(\oplus_3) = 9$; witnesses &
Lean 4 kernel (\texttt{decide}) for witnesses + DP structural framework; the finite exclusion
$\mathrm{par3} \notin D_8$ (hence $\tree \ge 9$) by \texttt{native\_decide} (compiler-trusted); two
implementations of the same recurrence agree \\
\addlinespace
$\gapop(\oplus_3) \ge 2$ &
Analytic (Khrapchenko + six-gate circuit), independent of the tree DP and the DRAT chain; and
fully citational via Schnorr + Khrapchenko \\
\addlinespace
$\gapop(\oplus_3) = 3$ (exact) &
Certified here, not cited: the DRAT chain for $\opt = 6$, and two independent dynamic programs
plus the Lean development for $\tree = 9$. No determination of $\tree(\oplus_3)$ appears in the
refereed literature (\S\ref{sec:refutation}); the value agrees with the exhaustive enumeration
of~\cite{cox2011} \\
\addlinespace
$\gapop(\oplus_4) = 6$ (exact, census maximum) &
Fully citational: Schnorr~\cite{schnorr1976} for $\opt = 9$ and Tarui~\cite{tarui2010} for
$\tree = 15$, since $4 = 2^2$ is where Khrapchenko's bound is attained \\
\addlinespace
$\opt(\oplus_3) = 6$, $\opt(\oplus_4) = 9$ (external agreement) &
Classical: the AIG basis is $U_2$ and Schnorr's theorem gives the $U_2$ complexity of $\oplus_n$ as
exactly $3n-3$~\cite{schnorr1976} --- an independent check of the encoder at two points \\
\addlinespace
AIG gap census $n = 4$ &
Two implementations of the tree recurrence agree on all 65{,}536 cells; the $\opt$ side is the
author's companion catalogue~\cite{krinkin-catalog}, re-derived here class by class (222/222) --- its UNSAT
side is gated but not DRAT-certified (declared) \\
\addlinespace
XAG census $n \le 4$ &
Gated encoder (independent exhaustive enumeration at $n = 3$; formula encoder matches DP on all 222
$n = 4$ classes); SAT models simulation-verified; UNSAT side not proof-certified (declared) \\
\addlinespace
XAG search $n = 5$ &
Encoder gated at $n \le 4$; 494 instances (all gap1 + a gap0 sample) re-decoded and
simulation-verified; sample, not exhaustive (declared) \\
\addlinespace
MIG census $n = 4$ &
Same gating + exact match with an independently published size distribution (aggregate) \\
\bottomrule
\end{tabularx}
\end{table}

\noindent
The two $\tree$ implementations are independent \emph{implementations of the same Bellman recurrence}
(layered vs.\ fixed-point) --- strong bug detection; the semantic bridge (that the recurrence
computes formula size) is the Lean completeness lemma. The four-model-family adversarial review is
provenance, not part of the verification chain.

\section{Related work and novelty}

\begin{itemize}
  \item \textbf{Krinkin, arXiv:2603.08033}~\cite{krinkin-unitgap} is the refuted paper;
  \textbf{arXiv:2603.09379}~\cite{krinkin-catalog} (an NPN-4 AIG catalog whose two open entries we
  closed) is a companion note. Both were communicated to the author, who
  confirmed the refutation independently; the exact values closing the catalogue's two open
  entries, and a proof for the unit-gap paper's Theorem~3, were contributed back
  (\texttt{krinkin/bounds\#2}, \texttt{krinkin/unit-gap\#2}).

  \item \textbf{MIG:} the optimal $n = 4$ sizes were previously published (Soeken et al., DATE
  2016~\cite{soeken2016}; TCAD 2017~\cite{soeken2017}); our census is replication, claimed only as
  an aggregate cross-check.

  \item \textbf{XAG total-gate count} is essentially the classical combinational complexity over
  $B_2$, tabulated for small $n$ by Knuth (TAOCP 7.1.2)~\cite{knuth-taocp4a}; our maximum of 7
  agrees. We claim \textbf{no novelty for the XAG numbers themselves}.

  \item \textbf{Both sizes of parity are known, and neither is ours.} Schnorr~\cite{schnorr1976}
  determined the $U_2$ circuit complexity of $\oplus_n$ exactly, $3n-3$. On the formula side
  Khrapchenko's $n^2$ lower bound is classical and, for $n = 2^k$, attained~\cite{tarui2010}.
  Since the AIG basis with free inversions is $U_2$ (\S\ref{sec:prelim}), that fixes
  $\opt(\oplus_3) = 6$ and $\tree(\oplus_3) \ge 8$ --- already enough to refute Theorem~2 --- and
  fixes the $n = 4$ values exactly ($\opt = 9$, $\tree = 15$, gap 6, the census maximum). For
  $\tree(\oplus_3) = 9$ we found no determination in the refereed literature, but the value is not
  new either: the exhaustive enumeration of Cox and Healy~\cite{cox2011} reports parity as a
  hardest function at both $n = 3$ and $n = 4$, needing 9 and 15 AND/OR operators, and exhibits
  the optimal $n = 3$ formula. Both agree with the values computed here. We claim \textbf{no
  novelty for any of these sizes}; they are used as an external check and to make the refutation
  citational. What
  we did not find in the literature is the two placed side by side: the circuit bound and the
  formula bound sit in different corners --- gate-elimination and rank methods respectively --- and
  reading them together requires only the basis identification above, which is why the
  contradiction with Theorem~2 could sit unnoticed in published results for two decades.

  \item \textbf{Exhaustive catalogs of small functions} --- exact synthesis over NPN-4 classes
  (Soeken et al.~\cite{soeken2018}; Haaswijk~\cite{haaswijk-thesis}), SAT-based circuit
  search~\cite{kojevnikov2009}, and ESOP exact synthesis~\cite{riener2018} --- tabulate optimal
  \emph{circuits} (or, for ESOP, optimal formulas alone). We did not find a catalog reporting both
  measures for the same functions. The closest related line is the study of \textbf{fanout-free
  (read-once) functions}~\cite{hayes1975}, which concerns formulas in which each \emph{variable}
  occurs once. That is strictly stronger than $\gapop(f) = 0$, which only asks that some optimal
  circuit have fan-out one at every \emph{gate}, variables being free to reappear as leaves:
  $x_1 \oplus x_2$ has $\gapop = 0$ in AIG ($\opt = \tree = 3$) and is not read-once. Neither line
  tabulates the gap.

  \item \textbf{Contributions reported here:} (i) the refutation of Theorem~2 and Corollary~6, with
  machine-checked certificates and a Lean formalization, and the localization of the error to the
  paper's own
  two incompatible characterizations in \S2 and \S3; and (ii) an \emph{explicitly tabulated}
  gap-by-basis comparison ($\tree - \opt$ per NPN class) in AIG and XAG at $n \le 4$, with a
  directed $n = 5$ XAG search. The $\tree - \opt$ difference is a subtraction of columns each of
  which is classical, so (ii) is a contribution of organized data and the base-dependence
  observation, not of a new measure; it is likely implicit in any catalog publishing both sizes. We
  searched Knuth's $B_2$ data, the MIG/AIG and NPN-4 exact-synthesis catalogs, the ESOP and
  read-once/fanout-free literature, and the standard EDA benchmark collections (IWLS, the EPFL
  combinational suite, LGSynth) without finding the
  gap-by-basis comparison stated explicitly, but treat (ii) as ``apparently unreported'' pending a
  specialist check.
\end{itemize}

\section{Open questions}

\begin{enumerate}
  \item Does $\gapx$ remain in $\{0,1\}$ for all of $n = 5$ (the directed search found no
  separator among 25{,}169 gap-decided functions) and for larger $n$? Parity no longer separates in
  XAG; other separators are not excluded.
  \item Is there a theorem behind the empirical XAG unit gap at small $n$ --- e.g.\ a normal-form
  argument bounding sharing once linear structure is factored out?
  \item An efficient exact DP for $\treemig$ (the deferred MIG gap column).
  \item A proof, or a counterexample, for the conditional statement of Krinkin's Theorem~7 under
  the standard definition. What $\oplus_3$ removes is the published proof, which routes through
  Corollary~6; the statement itself is left open.
\end{enumerate}

\section*{Acknowledgements}
I thank Kirill Krinkin. Told that Theorem~2 of~\cite{krinkin-unitgap} was false, he checked the
counterexample himself --- by Khrapchenko's bound, independently of anything presented here ---
confirmed it the same day, and said so plainly. He has stated that he will withdraw the preprint.
His reading of the damage is the one adopted in \S\ref{sec:refutation}, including the distinction
that is easiest to get wrong: Corollary~6 falls with Theorem~2, and Theorem~7 loses its proof but
is not itself refuted, while Theorems 3 and 4 stand. He raised no question of priority.

Conceding a refutation quickly and in full is harder than it looks, and rarer than it should be.
His companion catalogue~\cite{krinkin-catalog} also supplied the $n = 4$ optima that the census
rests on; its two open entries were closed in the course of this work.

\section*{Methods note}
The analysis, code and certificates reported here were produced by an AI system operating under
the direction of the author, who selected the target, set the verification standard and is
responsible for all claims. Correctness does not rest on that process: the circuit lower bounds
carrying the refutation --- $\opt(\oplus_3) \ge 6$ and the shared-child bound --- are DRAT proofs
checked by \texttt{drat-trim} (the census UNSAT sides are gated but not proof-certified, as
\S\ref{sec:verification} declares), the formula-size counterexample is checked by the
Lean~4 kernel except where \texttt{native\_decide} is declared, and the analytic fail-safe
(Khrapchenko) is independent of both. Before the results were communicated externally, the package
was reviewed adversarially by four independent language-model families; that review is provenance,
not verification.

\nocite{khrapchenko1971,kojevnikov2009,tools}
\bibliographystyle{plain}
\bibliography{refs}

\appendix

\section{Corrected proof of Theorem 3 (Threshold)}
\label{app:thm3}

\begin{claim}
If $f$ has $n$ essential variables and $\gapop(f) = 1$, then $\opt(f) = m \ge n$.
\end{claim}

\begin{proof}
Since $\gapop(f) = 1$, take a size-optimal single-output AIG $C$ for $f$. $C$ is \textbf{not a
tree}: if it were, $\tree(f) \le |C| = \opt(f)$, forcing $\gapop(f) = 0$, contrary to hypothesis. So
some non-output gate of $C$ has fan-out $\ge 2$. Now count input incidences. Every gate has two
inputs, so the total number of input slots is $2m$; partition them into $I_{\mathrm{var}}$ slots
reading a primary input or the constant, and $I_{\mathrm{gate}}$ slots reading another gate, so
$2m = I_{\mathrm{var}} + I_{\mathrm{gate}}$.

\begin{itemize}
  \item \emph{Essential variables:} each of the $n$ essential variables must feed at least one slot
  (else $f$ would not depend on it), so $I_{\mathrm{var}} \ge n$.
  \item \emph{Gate usage:} each of the $m - 1$ non-output gates feeds at least one later slot
  (normalized circuit: no unused gate), contributing $\ge m - 1$; the fan-out-$\ge 2$ gate
  contributes at least one further incidence, so $I_{\mathrm{gate}} \ge m$.
\end{itemize}

Hence $2m = I_{\mathrm{var}} + I_{\mathrm{gate}} \ge n + m$, i.e.\ $m \ge n$.
\end{proof}

\noindent
(The bound holds for any non-tree optimal single-output circuit, subsuming the $\gapop = 1$ case.
The published proof's $|S| \ge k-1$ step applies gate elimination to the cone $S$ \emph{excluding}
the shared gate $g$; when $g$ is at input level the $k$ essential variables of $S \cup \{g\}$ are
not all covered by $S$, and the count is short by one. The incidence argument avoids the case
split.)

\section{Encoders, gates, and the parity-5 false positive}
\label{app:encoders}

The repository documents the AIG, XAG and multi-output encoders together with the validation
gates each had to pass before use. For the XAG formula encoder these were: the complete $n = 3$
check ($\treex = \optx$ on all 256 functions), and \texttt{tree\_via\_sat} matching the independent
DP on all 222 classes at $n = 4$, including the four with $\gapx = 1$.

\paragraph{The parity-5 false positive.} The first $n = 5$ search implementation tested
formula-satisfiability at exactly $\opt + 1$ gates in isolation and declared a separator when that
single instance was UNSAT. This is unsound for a normalized exact-size encoder: UNSAT at exactly
$\opt + 1$ means ``no formula of \emph{exactly} $\opt + 1$ gates'', not $\tree \ge \opt + 2$ (a
formula of a different size may exist). The bug surfaced on parity-5 (\texttt{0x69969669}), a linear
function with $\gapx = 0$, which the isolated test misreported as a separator. The fix searches
\textbf{ascending from $\opt$} ($\tree$ is the first $k \ge \opt$ with a satisfiable formula
encoding); parity-5 then correctly returns gap 0. All results in \S\ref{sec:basedep} use the fixed procedure.

\section{Reproducibility}
\label{app:repro}

\begin{sloppypar}
Every artifact named below is public at
\url{https://github.com/totobusnello/unit-gap-certificates}; the repository root carries a
table mapping this paper's principal claims to the files that back them. Re-checking that the archived
CNFs are unsatisfiable requires \texttt{drat-trim} alone --- not the encoders, not the solver,
and nothing else in the repository. What a DRAT proof does \emph{not} certify is that a CNF
faithfully encodes the question asked of it; that is a separate obligation, discharged here by
cross-checking the encoder against independent exhaustive enumeration
(\S\ref{sec:verification}). The recorded transcript carries the SHA-256 of every CNF and DRAT
file, so a reader can confirm they are checking the same bytes.
\end{sloppypar}

\begin{sloppypar}
kissat 4.0.4; \texttt{drat-trim} (marijnheule @ commit \texttt{2e3b2dc}); Lean 4.31.0. Versioned
artifacts: encoders, gates, censuses (AIG \texttt{npn4\_gap.csv}, XAG \texttt{npn4\_xag\_gap.csv}
with per-class \texttt{opt\_xag}/\texttt{tree\_xag}/\texttt{gap\_xag}, MIG); CNFs and the
$k = 1,\dots,5$ DRAT proofs for $\opt(\oplus_3)$ and $k = 1,\dots,3$ for the representative child;
the \texttt{drat-trim} transcript \texttt{certificates/verify\_lowk.log} (\texttt{s VERIFIED}, exit code,
and CNF/DRAT SHA-256 for all eight proofs); the $n = 5$ search CSVs (raw per-worker and globally
deduplicated) and the 494/494 re-verification log; and the Lean sources. Both implementations of
the $\tree$ recurrence are included --- the layered dynamic program and the independent global
Bellman fixed point --- together with the transcript recording their agreement on all
65{,}536 cells. The rerun of the source catalogue's own checker is under
\texttt{verification/}; its two inputs and the checker itself are \emph{not} redistributed
(that repository carries no licence), and the SHA-256 needed to confirm them are. The $n = 5$
search seeds are deterministic per worker; the bench is a 2-vCPU VPS.
\end{sloppypar}

\end{document}
